\documentclass[12pt, a4paper]{article}

% ── Packages ──────────────────────────────────────────────────────────────────
\usepackage{amsmath, amssymb, amsthm}
\usepackage{geometry}
\usepackage{booktabs}
\usepackage{array}
\usepackage{hyperref}
\usepackage{xcolor}
\usepackage{tcolorbox}
\usepackage{enumitem}
\usepackage{parskip}

\geometry{margin=2.5cm}

% ══════════════════════════════════════════════════════════════════════════════
% ── Student-mode switch ───────────────────────────────────────────────────────
%
%   Compile with: \def\studentflag{1}\input{...}  → student version
%   Compile with: \input{...}                      → instructor version
%
\newif\ifstudentmode
\studentmodefalse       % default: instructor mode

\ifx\studentflag\undefined
  % no flag passed — keep studentmodefalse
\else
  \studentmodetrue
\fi
% ══════════════════════════════════════════════════════════════════════════════

% ── Student-mode macros ───────────────────────────────────────────────────────
%
%  \sol[<width>]{<math content>}   — blank inside a math environment
%  \soleq[<height>]{<math content>} — blank replacing an entire display equation
%
\newcommand{\sol}[2][6cm]{%
  \ifstudentmode
    \underline{\hspace{#1}}%
  \else
    {#2}%
  \fi
}

\newcommand{\soleq}[2][3cm]{%
  \ifstudentmode
    \par\vspace{4pt}%
    \noindent\rule{\linewidth}{0.5pt}\\[#1]%
    \noindent\rule{\linewidth}{0.5pt}%
    \par\vspace{4pt}%
  \else
    \[#2\]%
  \fi
}
% ──────────────────────────────────────────────────────────────────────────────

% ── Colors & boxes ────────────────────────────────────────────────────────────
\definecolor{mainblue}{RGB}{30, 100, 180}
\definecolor{boxbg}{RGB}{235, 243, 255}
\definecolor{notebg}{RGB}{255, 248, 230}

\tcbuselibrary{skins, breakable}

\newtcolorbox{resultbox}{
  colback=boxbg, colframe=mainblue,
  boxrule=1.2pt, arc=4pt,
  title={Result},
  fonttitle=\bfseries\color{white},
  coltitle=white, attach boxed title to top left={yshift=-2mm, xshift=4mm},
  boxed title style={colback=mainblue, arc=3pt}
}

\newtcolorbox{notebox}{
  colback=notebg, colframe=orange!70!black,
  boxrule=1pt, arc=4pt,
  left=4pt, right=4pt
}

% ── Theorem environments ───────────────────────────────────────────────────────
\theoremstyle{definition}
\newtheorem{example}{Example}[section]

% ── Math macros ───────────────────────────────────────────────────────────────
\newcommand{\bth}{\boldsymbol{\theta}}
\newcommand{\bx}{\mathbf{x}}
\newcommand{\bX}{\mathbf{X}}
\newcommand{\by}{\mathbf{y}}
\newcommand{\bw}{\mathbf{w}}
\newcommand{\yhat}{\hat{y}}
\newcommand{\sigmoid}{\sigma}

% ── Document ──────────────────────────────────────────────────────────────────
\begin{document}

\begin{center}
  {\LARGE \textbf{Section 1.2: Logistic Regression}} \\[0.5em]
  {\large AI Class — Lecture Notes}%
  \ifstudentmode
    \quad {\normalsize\color{gray}[Student Copy]}%
  \fi
\end{center}

\vspace{1em}
\hrule
\vspace{1.5em}

% ══════════════════════════════════════════════════════════════════════════════
\section{From Regression to Classification}

In Section~1.1 we predicted a \emph{continuous} output $y \in \mathbb{R}$.
Now suppose the output is a \textbf{binary label}: $y \in \{0, 1\}$.
For example, we might ask:

\begin{center}
  \renewcommand{\arraystretch}{1.2}
  \begin{tabular}{ll}
    \toprule
    \textbf{Input $x$} & \textbf{Label $y$}        \\
    \midrule
    Tumor size (cm)    & 0 = benign, 1 = malignant \\
    Hours studied      & 0 = fail, 1 = pass        \\
    Email features     & 0 = not spam, 1 = spam    \\
    \bottomrule
  \end{tabular}
\end{center}

\textbf{Why not use linear regression?}
If we fit $\yhat = wx + b$ directly and threshold at $0.5$, the model is
not constrained: it can predict values far below 0 or above 1, which are
meaningless as probabilities and sensitive to outliers.
What we really want is a model that outputs a \emph{probability}:
\[
  \hat{p} = P(y = 1 \mid x) \in (0, 1).
\]

% ══════════════════════════════════════════════════════════════════════════════
\section{The Sigmoid Function: Motivation via Log-Odds}

A natural way to build a probability model is to assume that the
\textbf{log-odds} (logit) of the event is a linear function of the input.

The \textbf{odds} of an event with probability $p$ are defined as:
\[
  \text{odds} = \frac{p}{1 - p}.
\]

We assume the \textbf{log-odds} are linear:

\begin{equation}
  \log \frac{p}{1-p} = wx + b.
  \label{eq:logodds}
\end{equation}

This is a reasonable model: the log-odds can range over all of $\mathbb{R}$,
which matches a linear function perfectly.

\subsection{Solving for \texorpdfstring{$p$}{p}}

Exponentiating both sides of \eqref{eq:logodds} and solving for $p$:

\begin{align*}
  \frac{p}{1-p}     & = e^{wx+b}        \\[4pt]
  p                 & = (1-p)\,e^{wx+b} \\[4pt]
  p\,(1 + e^{wx+b}) & = e^{wx+b}
\end{align*}

% ── \sol applied: student completes the derivation ────────────────────────────
\begin{resultbox}
  \[
    \hat{p} \;=\; \sol[9cm]{\frac{1}{1 + e^{-(wx+b)}} \;=\; \sigmoid(wx+b)}
  \]
\end{resultbox}
% ─────────────────────────────────────────────────────────────────────────────

This function $\sigmoid(z) = \dfrac{1}{1+e^{-z}}$ is called the
\textbf{sigmoid} (or logistic) function.

\begin{notebox}
  \textbf{Key properties of $\sigmoid(z)$.}
  \begin{itemize}[nosep]
    \item Output is always in $(0, 1)$ — valid as a probability.
    \item $\sigmoid(0) = 0.5$: the decision boundary is where $wx + b = 0$.
    \item As $z \to +\infty$, $\sigmoid(z) \to 1$; as $z \to -\infty$,
          $\sigmoid(z) \to 0$.
    \item Its derivative has the elegant form
          $\sigmoid'(z) = \sigmoid(z)\bigl(1 - \sigmoid(z)\bigr)$
          (see Appendix~\ref{app:sigmoid-deriv}).
  \end{itemize}
\end{notebox}

% ══════════════════════════════════════════════════════════════════════════════
\section{The Logistic Regression Model}

For $d$ features, the model generalises directly:

\begin{equation}
  \hat{p}^{(i)} = \sigmoid\!\left(\bth^{\top}\bx^{(i)}\right)
  = \frac{1}{1 + e^{-\bth^{\top}\bx^{(i)}}},
  \label{eq:logreg}
\end{equation}

where $\bth = [b,\, w_1,\, \ldots,\, w_d]^{\top}$ and $\bx^{(i)}$ includes a
leading 1 for the bias (same design-matrix convention as Section~1.1).

The predicted class is:
\[
  \hat{y}^{(i)} =
  \begin{cases} 1 & \text{if } \hat{p}^{(i)} \ge 0.5 \\ 0 & \text{otherwise.} \end{cases}
\]

% ══════════════════════════════════════════════════════════════════════════════
\section{Cost Function: Cross-Entropy via MLE}

\subsection{Setting Up the Likelihood}

We model each label $y^{(i)}$ as a Bernoulli random variable with success
probability $\hat{p}^{(i)}$:
\[
  P\!\left(y^{(i)} \mid \bx^{(i)};\,\bth\right)
  = \left(\hat{p}^{(i)}\right)^{y^{(i)}}
  \left(1 - \hat{p}^{(i)}\right)^{1 - y^{(i)}}.
\]

Assuming the $n$ observations are independent, the joint likelihood is:
\[
  \mathcal{L}(\bth)
  = \prod_{i=1}^{n}
  \left(\hat{p}^{(i)}\right)^{y^{(i)}}
  \left(1 - \hat{p}^{(i)}\right)^{1 - y^{(i)}}.
\]

\subsection{From MLE to the Cross-Entropy Loss}

Maximising $\mathcal{L}$ is equivalent to maximising the log-likelihood
$\ell = \log \mathcal{L}$ (log is monotone):

% ── \sol applied: student expands the log-likelihood ─────────────────────────
\[
  \ell(\bth)
  = \sol[10cm]{\sum_{i=1}^{n}
    \Bigl[y^{(i)} \log \hat{p}^{(i)}
      + (1 - y^{(i)}) \log\!\left(1 - \hat{p}^{(i)}\right)\Bigr].}
\]
% ─────────────────────────────────────────────────────────────────────────────

Minimising the \textbf{negative} log-likelihood (and averaging over $n$) gives
the \textbf{Binary Cross-Entropy (BCE)} loss:

\begin{resultbox}
  \[
    J(\bth)
    = \sol[10cm]{-\frac{1}{n}\sum_{i=1}^{n}
      \Bigl[y^{(i)} \log \hat{p}^{(i)}
        + (1 - y^{(i)}) \log\!\left(1 - \hat{p}^{(i)}\right)\Bigr]}
  \]
\end{resultbox}

\begin{notebox}
  \textbf{Why cross-entropy?}
  When $y^{(i)} = 1$, only the first term $-\log\hat{p}^{(i)}$ matters:
  we are penalised heavily if $\hat{p}^{(i)}$ is close to 0 (confident and
  wrong).  When $y^{(i)} = 0$, only the second term $-\log(1-\hat{p}^{(i)})$
  matters.  The two terms switch on and off depending on the true label.
\end{notebox}

% ══════════════════════════════════════════════════════════════════════════════
\section{The Gradient and Why We Need Gradient Descent}

Unlike MSE in linear regression, there is \textbf{no closed-form solution}
for the minimiser of $J(\bth)$.  The cross-entropy loss is convex (it has a
unique global minimum), but we cannot invert the normal equations here because
the sigmoid is non-linear.

Taking the gradient with respect to $\bth$ (using the chain rule and the
identity $\sigmoid'(z) = \sigmoid(z)(1-\sigmoid(z))$):

\begin{resultbox}
  \[
    \nabla_{\bth}\,J
    = \sol[8cm]{\frac{1}{n}\sum_{i=1}^{n}
      \Bigl(\hat{p}^{(i)} - y^{(i)}\Bigr)\,\bx^{(i)}}
  \]
\end{resultbox}

\begin{notebox}
  \textbf{Remarkably clean form.}
  The gradient has exactly the same structure as in linear regression —
  it is a weighted sum of the input vectors, where the weights are the
  \emph{prediction errors} $(\hat{p}^{(i)} - y^{(i)})$.
  The sigmoid ``absorbed'' all the complexity through the chain rule.
\end{notebox}

Because there is no closed form, we use \textbf{gradient descent}: starting
from an initial $\bth^{(0)}$, we repeatedly update
\[
  \bth \;\leftarrow\; \bth - \alpha\,\nabla_{\bth}\,J,
\]
where $\alpha > 0$ is the \textbf{learning rate}.
Gradient descent will be covered in detail in the next section.

% ══════════════════════════════════════════════════════════════════════════════
\section{Summary}

\begin{center}
  \renewcommand{\arraystretch}{1.5}
  \begin{tabular}{ll}
    \toprule
    \textbf{Model}           & $\hat{p} = \sigmoid(\bth^{\top}\bx)$                                        \\
    \textbf{Loss function}   & BCE: $J = -\dfrac{1}{n}\sum\bigl[y\log\hat{p} + (1-y)\log(1-\hat{p})\bigr]$ \\
    \textbf{Gradient}        & $\nabla_{\bth}J = \dfrac{1}{n}\sum(\hat{p}^{(i)} - y^{(i)})\,\bx^{(i)}$     \\
    \textbf{Solution method} & Gradient descent (no closed form)                                           \\
    \textbf{Works well when} & The decision boundary is (approximately) linear                             \\
    \bottomrule
  \end{tabular}
\end{center}

\bigskip
\textbf{Up next (Section~1.3):} What if no linear boundary can separate the
classes?  We revisit the XOR dataset to show that both linear regression and
logistic regression fail — motivating the need for non-linearity and neural
networks.

% ══════════════════════════════════════════════════════════════════════════════
\appendix
\section{Derivative of the Sigmoid Function}
\label{app:sigmoid-deriv}

Let $\sigmoid(z) = (1 + e^{-z})^{-1}$.  By the chain rule:

\begin{align*}
  \sigmoid'(z)
   & = -\frac{-e^{-z}}{(1+e^{-z})^2}
  = \frac{e^{-z}}{(1+e^{-z})^2}.
\end{align*}

Factor out $\dfrac{1}{1+e^{-z}}$ from the numerator and denominator:

\[
  \sigmoid'(z)
  = \frac{1}{1+e^{-z}} \cdot \frac{e^{-z}}{1+e^{-z}}
  = \sigmoid(z) \cdot \frac{1+e^{-z}-1}{1+e^{-z}}
  = \sigmoid(z)\bigl(1 - \sigmoid(z)\bigr).
\]

\begin{resultbox}
  \[
    \sigmoid'(z) = \sigmoid(z)\bigl(1 - \sigmoid(z)\bigr)
  \]
\end{resultbox}

This compact form makes the chain-rule calculations in gradient descent
particularly clean.

\end{document}
